\documentclass[11pt]{article}
\usepackage[margin=1.1in]{geometry}
\usepackage{amsmath, amssymb}
\usepackage{booktabs}
\usepackage{xcolor}
\usepackage{algorithm}
\usepackage{algpseudocode}
\usepackage{enumitem}

\title{TauBots: graded transparency over the zoo\\
\large Design, experimental setup, and first results}
\author{Colomban Duclaux}
\date{\today}

% Bot names, matching the other notes.
\newcommand{\zoobot}[1]{\texttt{#1}}
\newcommand{\dupoc}{\zoobot{DupocBot}}
\newcommand{\cupod}{\zoobot{CupodBot}}
\newcommand{\cimcic}{\zoobot{CIMCIC}}
\newcommand{\dimcid}{\zoobot{DIMCID}}
\newcommand{\Tau}{\mathrm{Tau}}

\begin{document}
\maketitle

\section{Introduction}
In these notes, we describe an experiment on graded transparency based on the TauBots framework. The headline finding is that transparency's effect on cooperation depends on the agents' caution
threshold \emph{in sign}: cautious agents need transparency to cooperate at
all, trusting agents cooperate \emph{more} as the signal degrades, and between
the two lies a band where the tournament barely moves. On the larger zoo we
also observe a phenomenon we call \emph{floor healing}: blur restores
cooperation between prover bots that full transparency, combined with bounded
proof budgets, had destroyed.

\section{Design choices}

\subsection{Why partial transparency, and what constrains the design}

Two constraints shaped everything that follows. First, we must \textbf{reuse
the existing zoo and outcome matrix}. Prover bots (\dupoc{} and company) need
\emph{exact syntax} to prove anything --- a guard like ``the opponent
provably cooperates'' quantifies over a concrete term, and a corrupted or
noised source is simply useless to them. Whatever ``blur'' means, it cannot
mean perturbing programs. Second, we want \textbf{graded intensity}: the
formalism should let us speak meaningfully of 20\%, 40\%, 60\% transparency,
so that cooperation can be plotted \emph{against} transparency rather than
merely compared at the two endpoints.

The first constraint essentially forces the design. If hypotheses must be
exact programs, then uncertainty can only live in \emph{which} program the
opponent is. A signal is therefore a probability distribution over the zoo:
candidates $B_1, \dots, B_n$ with weights $p_1, \dots, p_n$ summing to one.
Every hypothesis is a real, fully-transparent zoo member backed by a proven
matrix cell; the blur lives entirely in the weights.

\subsection{Defining the tau lift: a false start and the fix}

Given a signal, how should an agent $A$ act on it? After some research, the
definition one first writes down is a threshold vote over the hypotheses:

\begin{algorithm}[h]
\caption{First attempt (rejected): tau hypotheses in the signal}
\begin{algorithmic}
\Procedure{$\Tau A(\alpha)$}{$\Tau B_1,\dots,\Tau B_n,\; p_1,\dots,p_n$}
  \State $\mathit{sumCs} \gets 0$
  \For{$i = 1, \dots, n$}
    \If{$\mathrm{outcome}(A(\Tau B_i)) = C$}
      \State $\mathit{sumCs} \mathrel{+}= p_i$
    \EndIf
  \EndFor
  \If{$\mathit{sumCs} \ge \alpha$} \Return $C$ \Else{} \Return $D$ \EndIf
\EndProcedure
\end{algorithmic}
\end{algorithm}

The flaw is hiding in the argument list: the hypotheses are themselves
\emph{tau} agents. Each $\Tau B_i$'s play depends on a signal about
\emph{its} opponent, whose play depends on a further signal, and the
definition recurses with no base case --- one must remove the inner
$\Tau B \to B$, because otherwise the recursion is infinite. This is not a
cosmetic problem: unlike the self-referential loops that L\"obian
cooperation resolves, this regress has no provability structure to anchor
it, so there is no bounded-L\"ob rescue. It is an ungrounded, potentially
bistable fixpoint, and we reject the definition outright.

The fix is to ground the hypotheses at the \emph{base} zoo:

\begin{algorithm}[h]
\caption{Final version (Def.~3): base hypotheses, self probe}
\begin{algorithmic}
\Procedure{$\Tau A(\alpha)$}{$B_1,\dots,B_n,\; p_1,\dots,p_n$}
  \State $\mathit{sumCs} \gets 0$
  \For{$i = 1, \dots, n$}
    \If{$\mathrm{outcome}(A(B_i)) = C$}
      \Comment{$A$'s \emph{own} action in the proven cell $\mathrm{outcome}(A, B_i)$}
      \State $\mathit{sumCs} \mathrel{+}= p_i$
    \EndIf
  \EndFor
  \If{$\mathit{sumCs} \ge \alpha$} \Return $C$ \Else{} \Return $D$ \EndIf
\EndProcedure
\end{algorithmic}
\end{algorithm}

The initial definition could in theory be used to define Tau agents, but it is much more complex to implement in practice, and it is not clear that it would be well-defined, which is why we stick with the final version for this experiment.

\subsection{Harsanyi type spaces -- Literature Review}

The construction is not ad hoc; it is a Harsanyi type space with the zoo as
the type set. Harsanyi's move for games of incomplete information was to
compress the infinite belief hierarchy into a \emph{type}, each player
holding a probability distribution over the opponent's possible types. Our
signal is exactly such a belief, and the signal's concentration is the
information dial: a point mass is complete information --- Critch's regime
--- while the uniform distribution is no information, the classical opaque
PD.

We deviate from the full Harsanyi apparatus in three deliberate ways. The
hierarchy is truncated at level one: the zoo is level zero, and there are no
beliefs about beliefs. There is no common prior: signals are exogenous
parameters we sweep, so our outputs are phase diagrams rather than
equilibria. And our types are \emph{intensional} --- a type determines the
opponent's \emph{term}, not merely its payoffs, which is precisely what lets
behaviorally identical types be distinguished (or not) by the L\"obian
fragment. Note also that the tau lift is \emph{not} Bayesian best response:
it thresholds $A$'s own cooperation mass and thereby preserves $A$'s
decision procedure. The information structure is Harsanyi's; the decision
rule deliberately is not.

\subsection{Does the lift commute with the base game?}

A natural question is whether
$\mathrm{outcome}(A(\Tau B)) = \mathrm{outcome}(A(B))$ --- whether playing
against a lifted opponent is the same as playing against the base one. The
conjectured answer is a clean split along the zoo's most meaningful
boundary. For \textbf{sim-only} agents $A$ (\zoobot{TitForTatBot},
\zoobot{MirrorBot}, \dots) the equation should hold for sufficient fuel:
simulators only observe behavior, and a lifted $B$ under a point-mass
counterfactual signal behaves exactly like $B$. \emph{Blur is invisible to
behavioral bots.} For \textbf{proof-search} agents $A$ (\dupoc,
\zoobot{PrudentBot}, \dots) the equation \emph{fails}: their guards are
provability statements about the opponent's literal term, and
$\Tau B$'s term is not $B$'s term. \emph{Blur is detectable precisely by
the bots that read source code.} If proven, this split would both justify
the construction --- the new phenomena live exactly where Critch's
machinery lives --- and cut the experimental workload, since behavioral
base-vs-tau cells reduce to existing matrix cells.

As a general theorem the equation is simply false: open-source game theory
is intensional, proof costs scale with term size, and the mutual-L\"ob
fixpoint shapes do not exist against tau-shaped terms. One could stipulate
the equation as a \emph{definition} of how tau agents play --- but that
stipulation \emph{is} Def.~3, which is why we write
$\mathrm{outcome}(A, B_i)$ and move on.

\section{Experimental setup}

The tau layer is pure matrix arithmetic over the Lean-verified outcome
table, so the experiments run entirely in Python with the kernel-checked
matrix as ground truth. Two ingredients vary: the signal family
$\sigma$, and the sub-zoo the matrix is restricted to. We first define the
three signal families with their raw internal knobs, then explain how those
knobs are calibrated onto the single transparency scale $t$ that all the
results are plotted against.

\subsection{Three $\sigma$ families}

Except for the null control, each family has the form
$\sigma(B)[B'] \propto \exp(-d(B,B')/T)$ for a family-specific distance
$d$ and an internal temperature $T$ (a small $T$ concentrates the signal
on the true opponent; a large $T$ flattens it toward uniform).

The \textbf{behavioral} family takes $d$ to be the Hamming distance between
the two bots' own-action rows in the outcome matrix: close means ``plays
like''. This is not an arbitrary choice of kernel --- it is the exact Bayes
posterior (under a uniform prior) of an observer who watched the bot play
every zoo member through a noisy channel that flips each observation with
probability $p$, since $p^d(1-p)^{n-d} \propto \exp(-d \cdot
\ln\frac{1-p}{p})$ is a softmax in $d$. The behavioral channel is the
formalization of the classical ``watching it play'' story.

The \textbf{$\varepsilon$-uniform} family,
$\sigma = (1-\varepsilon)\,\delta_B + \varepsilon\,U$, is the null control.
It has no similarity structure at all: every bot is equally confusable with
every other, so identity leaks but resemblance plays no role. Because it is
identity-based it has no twin ceiling (below), and divergence between it
and a structured family at matched transparency isolates the effect of
confusion structure itself.

The \textbf{syntactic} family models an observer who partially \emph{reads
the source} rather than watching play --- degraded open-source game theory
in the most literal sense. Here $d$ is the normalized Zhang--Shasha tree
edit distance between parsed \texttt{Prog} terms: the minimum number of
node insertions, deletions and relabelings turning one abstract syntax tree
into the other, divided by the larger tree's node count so that the measure
is scale-free. The syntactic
channel's confusion structure \emph{inverts} the behavioral one ---
\dupoc{} and \cupod{} are syntactic near-twins (same tree, swapped action
leaves) with opposite dispositions, while \zoobot{CooperateBot} and
\zoobot{CupodTrollBot} are behavioral near-twins with wildly different
sources --- and globally the two distances correlate at only $+0.16$, so
the two channels are genuinely different objects, not noisy copies.

\subsection{The transparency dial}

The three families now on the table expose three incomparable knobs: two
softmax temperatures over different distances, and a mixing weight
$\varepsilon$. ``Temperature $2.7$ on the behavioral channel'' and
``$\varepsilon = 0.3$'' say nothing about how much information each signal
actually carries, so comparing families --- or even plotting one family on
a meaningful axis --- requires a common scale. That scale is the
transparency dial $t \in [0,1]$, and it \emph{is} the normalized mutual
information $I(B;\sigma(B))/H(B)$ of the signalling channel under a
uniform prior over the zoo: the fraction of the uncertainty about the
opponent's identity that the signal resolves, with $t = 1$ full
transparency and $t = 0$ opaque. Each family's internal knob is
bisection-inverted against this scale and never exposed: dialing
$t = 0.4$ yields, in every family, a channel whose \emph{measured}
transparency is 40\%. The payoff of this calibration discipline is
comparability: at matched $t$, every signal family carries the \emph{same
information rate}, so any difference between families at equal $t$
isolates the effect of confusion \emph{structure} --- which bots get
mistaken for which --- rather than differences in how much information
flows.

\subsection{Two sub-zoos}

The tau layer demands a \emph{totally proven} ordered submatrix, so zoos
are chosen as closed fragments of the outcome table rather than by fiat.

The \textbf{default zoo} (11 bots: \zoobot{CooperateBot}, \cupod,
\zoobot{CupodTrollBot}, \zoobot{DBot}, \zoobot{DefectBot}, \dupoc,
\zoobot{EBot}, \zoobot{GuardianBot}, \zoobot{OBot}, \zoobot{PrudentBot},
\zoobot{TitForTatBot}) is the clean instrument. It carries two stipulated (unproven) cells involving
\cupod, flagged as such by the tooling.

The \textbf{enlarged zoo} (16 bots) adds \zoobot{JustBot} and
\zoobot{LegibleBot} --- the behavioral twins deliberately excluded from
the default zoo --- together with \cimcic{} and \dimcid{} from the
search$\times$search frontier, and \zoobot{MirrorBot}. This purchase comes
with three assumptions worth stating plainly. Seven cells are stipulated
rather than proven (all on the search$\times$search frontier), so every
enlarged-zoo result is conditional on them. \zoobot{MirrorBot}'s self-play
is the proven-\texttt{none} fifth state $N$ --- mutual simulation never
terminates --- which we read \emph{pessimistically}: an $N$ hypothesis
contributes no cooperation mass. 
\section{Experimental results and analysis}

\subsection{The anchor theorem, and two instructive failures}

The construction's basic sanity check is the \emph{anchor}: at full
transparency the tau tournament should reproduce the base outcome matrix
exactly, since every signal is (numerically) a point mass and every
cooperation mass is $0$ or $1$. On the default zoo this holds on the nose:
at $t = 1$, $\alpha = 0.5$, zero cells differ from base. It is not clear whether this theorem would hold for Taubots as defined by the first attempt.

\subsection{The caution axis: three regimes and two knife edges}

Sweeping $\alpha$ at fixed transparency reveals the same three-regime
structure in both zoos, bracketed by two knife edges.

At the top, \textbf{$\alpha = 1$ collapses instantly}. Unanimity matches
the base matrix at $t = 1$, but at $t = 0.9$ mutual cooperation crashes to
$0.01$ (default) and $0.02$ (enlarged) --- exactly the unconditional
cooperators' self-cells, $1/121$ and $4/256$. A unanimous voter tolerates
zero blur: the moment any hypothesis mass falls on a defection-eliciting
bot, cooperation dies.

At the bottom, \textbf{$\alpha \to 0$ saturates --- with one telling
exception}. At full opacity and $\alpha = 0.1$, an agent cooperates unless
its \emph{base cooperation fraction} over the zoo falls below $\alpha$. In
the enlarged zoo only \zoobot{DefectBot} fails this test, giving
$\mathrm{CC} = (15/16)^2 = 0.88$. In the default zoo, however,
\zoobot{PrudentBot} fails it too: it cooperates with just one of eleven
members ($1/11 = 0.09 < 0.1$), so even a maximally trusting
tau-\zoobot{PrudentBot} defects under uniform blur, and the ceiling is
$(9/11)^2 = 0.67$. \zoobot{PrudentBot} is so selective that its lift is a
near-defector in the opaque limit --- prudence, pushed through the tau
construction, shades into misanthropy.

Between the edges lies \textbf{a dead zone where $\alpha$ is irrelevant}.
Only finitely many cooperation masses are achievable at a given $t$, and
between consecutive masses every $\alpha$ produces the identical
tournament. At $t = 0.5$ the gap is wide in the default zoo ---
$(0.41, 0.65)$, so the entire ``moderate caution'' band collapses to one
behavior --- and much narrower in the enlarged zoo, $(0.55, 0.65)$. With
158 distinct achievable masses against 98, the enlarged zoo simply offers
the caution dial more purchase. (A practical corollary for anyone
extending these sweeps: always step \emph{between} breakpoints; an
$\alpha$ sitting exactly on an achievable mass lets floating-point noise
decide plays.)

\subsection{The transparency axis: the effect on cooperation is
$\alpha$-dependent in sign}

The central result of the sweep is easy to state: \emph{whether
transparency helps cooperation depends on who is asking}. The same pattern
holds in both zoos and under all three signal families:

\begin{center}
\small
\begin{tabular}{lp{0.62\textwidth}}
\toprule
regime & effect of \emph{decreasing} transparency \\
\midrule
$\alpha \lesssim 0.4$ (trusting) & cooperation \textbf{rises} monotonically
  (default $0.33 \to 0.67$; enlarged $0.46 \to 0.88$ at $\alpha=0.1$) \\
$\alpha \approx 0.5$--$0.6$ & flat at the base rate over most of the range \\
$\alpha \gtrsim 0.7$ (cautious) & cooperation \textbf{falls} monotonically
  (both zoos $\to$ $\approx 0.01$--$0.06$) \\
\bottomrule
\end{tabular}
\end{center}

For cautious agents, transparency is \emph{necessary} for cooperation:
as the signal degrades, hypothesis mass leaks onto bots they would defect
against, and their threshold stops being met. For trusting agents,
transparency actively \emph{suppresses} cooperation: under blur they
extend cooperation to nearly anyone, including bots they would correctly
defect against under full information. And around $\alpha \approx 0.5$
the two forces roughly cancel, leaving a transparency-invariant band
anchored at the base rate. The question ``does transparency help
cooperation?'' has no unconditional answer; the honest unit of analysis
is the $(t, \alpha)$ phase diagram.

\subsection{What the extra cooperation is made of: suckers in one zoo,
healed provers in the other}

A rising mutual-cooperation rate does not say \emph{which} outcome it
displaces, so we decompose every tournament cell into mutual cooperation,
exploitation, and mutual defection. Here the two zoos tell genuinely
different stories at $\alpha = 0.5$.

In the \textbf{default zoo, blur feeds the defectors}. As transparency
falls toward $t = 0.3$, exploitation \emph{spikes} from $0.45$ to $0.53$
before mutual cooperation ever rises: the newly cooperative cells are
\cupod, \zoobot{EBot} and \zoobot{GuardianBot} extending cooperation
\emph{to} \zoobot{DefectBot}, which happily defects back. Meanwhile
\zoobot{PrudentBot} \emph{loses} its single cooperation (with \dupoc) ---
prudence collapses under blur. In this zoo, rising cooperation at low
transparency is substantially bots being suckered.

In the \textbf{enlarged zoo, blur heals prover--prover defection}. The
same sweep shows six cooperation gains, \emph{zero} losses, and none of
the gains directed at exploiters: the gainers are \cimcic{} (twice),
\cupod, \zoobot{CupodTrollBot}, \dupoc{} and \zoobot{JustBot} --- the
search fragment cooperating \emph{with each other}. Mutual defection
collapses from $0.20$ to $0.06$ while exploitation stays flat. The
mechanism deserves spelling out. Many of the enlarged zoo's $(D,D)$ cells
are not preference facts but \emph{proof-theoretic floor artifacts}: pairs
of provers that ``want'' to cooperate but cannot certify each other's
cooperation at matched budgets, because the certificate would have to pay
the very search floor it reasons about. Under blur, the hypothesis set
around such an opponent includes its cooperative behavioral neighbors;
the weighted vote crosses threshold; and the cooperation that proof
search could not deliver simply appears. \textbf{Ignorance restores
cooperation that the proof floor destroyed} --- to our knowledge the
cleanest example in this project of information reducing welfare through
the provability channel.

\subsection{Robustness: who needs transparency to be themselves}

For each bot we record the highest transparency at which its tau row first
departs from its base row --- a per-bot fragility measure (higher
threshold means more fragile). Three patterns emerge.

The constant bots (\zoobot{CooperateBot}, \zoobot{DefectBot},
\zoobot{LegibleBot}) never deviate at any $\alpha$: unconditional
behavior has nothing to lose to blur. Among the conditional bots,
fragility is \textbf{U-shaped in caution}: in the default zoo the most
robust regime is $\alpha = 0.5$, where every threshold sits at or below
$0.35$, while at $\alpha = 0.3$ and $\alpha = 0.8$ deviations begin as
early as $t = 0.65$--$0.75$. Mid-caution agents are the most faithful
lifts of their base selves --- yet another sense in which
$\alpha \approx 0.5$ is the construction's stable point. Moreover the
\emph{identity} of the fragile class rotates with $\alpha$: at high
caution the provers and probers break first (\dupoc, \zoobot{DBot},
\zoobot{EBot}, \cimcic, \zoobot{JustBot} at $0.75$--$0.80$), while at low
caution it is the guarded and trollish bots (\cupod, \zoobot{GuardianBot},
\zoobot{EBot}); \zoobot{OBot} is strikingly robust at low $\alpha$,
deviating only at full opacity. In the enlarged zoo, \dimcid{} and
\zoobot{MirrorBot} sit at threshold $1.00$ for every $\alpha$ --- these
are the structural twin and $N$ artifacts of \S 3.1, not blur
sensitivity, and should be read as such.

\subsection{Does confusion structure matter?}

Finally, the comparison the calibration discipline was built for: at
matched transparency, does it matter \emph{which} bots get confused? Above
$t \approx 0.5$, essentially not --- the three families agree almost
everywhere, so moderate blur acts through its information rate alone.
Below that, structure pays a small but consistent dividend: at
$\alpha = 0.5$ the behavioral channel exceeds the $\varepsilon$-uniform
null by $+0.017$ (default) and $+0.039$ (enlarged) in mutual cooperation.
Being confused with your behavioral neighbors is mildly
cooperation-favoring compared to being confused uniformly --- and twice
as much so in the enlarged zoo, whose confusion neighborhoods are more
cooperative. The syntactic channel tracks between the two. The effect is
real but second-order next to the $t$ and $\alpha$ axes.

\subsection{Head-to-head summary}

\begin{center}
\small
\begin{tabular}{lp{0.3\textwidth}p{0.3\textwidth}}
\toprule
 & default (11) & enlarged (16) \\
\midrule
base (CC / exploit / DD) & $0.33$ / $0.45$ / $0.22$ & $0.46$ / $0.34$ / $0.20$ \\
opaque limit, $\alpha=0.1$ & $0.67$ (excl.\ \zoobot{DefectBot}, \zoobot{PrudentBot}) & $0.88$ (excl.\ \zoobot{DefectBot}) \\
blur at $\alpha=0.5$ mainly\dots & feeds \zoobot{DefectBot} (exploitation spike at $t \approx 0.3$) & heals prover--prover: DD $0.20 \to 0.06$ \\
$\alpha$ dead zone at $t=0.5$ & wide, $(0.41, 0.65)$ & narrow, $(0.55, 0.65)$ \\
distinct $\alpha$ breakpoints & 98 & 158 \\
structure value at low $t$ & $+0.017$ & $+0.039$ \\
\bottomrule
\end{tabular}
\end{center}

The enlarged zoo is more cooperative at baseline, gains more from blur,
suffers less exploitation from it, and gives both dials more resolution.
But it buys all of this with seven stipulated cells, a broken anchor, and
two bots whose lifts are permanently distorted. Our working division of
labor is therefore: use the \textbf{default zoo} for any claim that leans
on the anchor theorem or on kernel-clean provenance, and the
\textbf{enlarged zoo} for the phenomena that need the search fragment ---
above all the floor-healing effect, which simply does not exist in the
default zoo's flip anatomy.

\subsection{Caveats and next steps}

All enlarged-zoo results are conditional on seven stipulated cells; the
tooling flags them, and a stipulation can never silently shadow a proven
cell (the loader fails loudly, which has already caught one stale entry in
practice). Grids are $t$ in steps of
$0.05$--$0.1$ and $\alpha$ in steps of $0.1$; the critical band around
$\alpha \approx 0.5$ deserves a finer sweep. And unless stated otherwise
the numbers above use the behavioral family --- whether the
prover-fragment story survives under the \emph{syntactic} channel, where
what leaks is code rather than conduct, is the most natural next
experiment.

\end{document}
